\section*{Chapitre \ref{ch:s1:volatility-targeting-and-risk-managed-portfolios} --- Ciblage de volatilité et portefeuilles pilotés par le risque}

\begin{solution}{exo:s1:volatility-targeting-and-risk-managed-portfolios:1}
$0.10/0.12 = 0.83$ et $0.10/0.40 = 0.25$.
\end{solution}

\begin{solution}{exo:s1:volatility-targeting-and-risk-managed-portfolios:2}
$0.002 \times (0.73 - 0.41)/0.004 = 0.16$~: environ 16~\% du volume de la journée (le 14.5~\% de la simulation retranche de l'exposition des fonds la dérive de prix de la journée).
\end{solution}

\begin{solution}{exo:s1:volatility-targeting-and-risk-managed-portfolios:3}
L'exposition varie avec le carré du rapport de la prévision~: une prévision supérieure de 10~\% réduit l'exposition d'environ 17~\% au lieu de 9~\%, de sorte que chaque variation de prévision est
une transaction plus importante.
\end{solution}

\begin{solution}{exo:s1:volatility-targeting-and-risk-managed-portfolios:4}
$0.06/0.16 = 0.375$. Un ratio de Sharpe sur dix ans a une erreur type d'environ $1/\sqrt{10} = 0.32$ (les échantillons vont de $-0.44$ à 1.18). Le gain moyen sur
vingt échantillons a une erreur type de $0.22/\sqrt{20} = 0.049$, donc 0.028 n'est pas significatif, même avec vingt décennies.
\end{solution}

\begin{solution}{exo:s1:volatility-targeting-and-risk-managed-portfolios:5}
Sur les actions, la volatilité augmente lorsque les cours baissent (l'effet de levier), si bien qu'elle est élevée après les pertes, et réduire l'exposition permet alors d'éviter la suite la plus
défavorable~; les obligations n'ont pas ce lien, si bien que leur volatilité dit peu de chose de leurs rendements futurs, et la mise à l'échelle n'ajoute que du bruit.
\end{solution}

\begin{solution}{exo:s1:volatility-targeting-and-risk-managed-portfolios:6}
Un pic de volatilité pousse les fonds à vendre~; leurs ventes font encore bouger les cours~; des mouvements plus amples relèvent les prévisions de volatilité des fonds~; des prévisions plus élevées
les poussent à vendre davantage. L'impact et la mémoire de la prévision déterminent jusqu'où la boucle se poursuit.
\end{solution}

\begin{solution}{exo:s1:volatility-targeting-and-risk-managed-portfolios:7}
Avant le pic, la volatilité réalisée du marché passe de 13.0~\% (sans impact) à 18.7~\%, et l'exposition des fonds tombe de 0.73 à 0.46~: leur propre
rééquilibrage quotidien, soit environ 1.9~\% du volume d'un jour moyen, fait bouger les cours, ce qui relève la volatilité qu'ils prévoient, ce qui les conduit à négocier davantage. Suffisamment
important, les fonds de ciblage engendrent une partie de la volatilité qu'ils visent.
\end{solution}

\begin{solution}{exo:s1:volatility-targeting-and-risk-managed-portfolios:8}
La stratégie implicite de la régression de spanning combine les portefeuilles pilotés et non pilotés, avec des pondérations estimées sur l'ensemble de l'échantillon~; elle n'est pas
investissable en temps réel. Il faut comparer directement le portefeuille piloté au portefeuille non piloté, hors échantillon, en intégrant les coûts, comme l'ont fait Cederburg et ses coauteurs.
\end{solution}

\begin{solution}{pb:s1:volatility-targeting-and-risk-managed-portfolios:1}
\begin{enumerate}
\item Un portefeuille ajusté par l'inverse de la volatilité prévue~; les flux nets de négociation des fonds qui suivent des règles mécaniques d'exposition.
\item Une prévision de variance, une règle d'ajustement, un cap d'effet de levier et une bande de rééquilibrage.
\item Les variations de volatilité ne sont pas compensées par des variations proportionnelles des rendements espérés~; prendre moins de risque lorsque la volatilité est élevée relève donc les ratios de Sharpe.
\item Sur 103 stratégies, les portefeuilles pilotés ne surpassent pas systématiquement les portefeuilles non pilotés hors échantillon.
\item Vingt échantillons de dix ans d'un marché GJR-GARCH à prime de risque constante, avec une exposition constante ou pilotée par la variance réalisée du mois précédent.
\item 0.49, 0.52 (inverse de la volatilité) et 0.48 (inverse de la variance)~; 10 et 8 échantillons sur 20 ont été améliorés.
\item Le gain (0.03) représente un dixième de l'erreur type sur une décennie (0.32).
\item Le gain ne tient que pour les actifs risqués dotés d'un effet de levier~; les classes synthétiques, qui n'en ont pas, varient dans les deux sens.
\item Des fonds détenant 0.2~\% de la valeur du marché et visant 10~\%, un marché échangeant 0.4~\% par jour, une baisse de 6~\% et vingt jours à 40~\% de volatilité~; tout est supposé.
\item 14.5~\% du volume du lendemain~; 20.5~\% du volume d'une journée sur vingt jours.
\item Il approfondit la baisse sur les vingt jours, de 16.6~\% à 18.7~\% (impact 0.1) ou 22.0~\% (0.3).
\item Cela découle de la volatilité récente du marché et des règles publiques des fonds.
\item \textbf{Résultat nommé.} Piloter le marché synthétique selon sa volatilité a relevé le ratio de Sharpe de 0.028 en moyenne sur vingt échantillons (meilleur dans
10 échantillons)~; après une baisse de 6~\%, les fonds de ciblage détenant 0.2~\% du marché ont vendu 14.5~\% du volume du lendemain.
\item Il réduit les queues de distribution et stabilise le risque.
\item Prévisions en temps réel et paramètres constants, comparaison directe avec le portefeuille non piloté, coûts, nombreuses périodes.
\item Négocier à contre-courant du flux de désendettement.
\item Le chapitre~19 visait la volatilité d'un portefeuille de tendance et constatait qu'il augmentait son effet de levier lorsque l'edge était faible~; ici la question est la même pour un marché.
\item Les actifs des fonds, leurs règles et leurs fenêtres de prévision, ainsi que le volume et la profondeur du marché à la clôture.
\item Leurs ventes alimentent leurs propres prévisions.
\item Un risque plus stable, acheté au prix de ventes dans les baisses, que le marché paie par des baisses plus profondes.
\end{enumerate}
\end{solution}

\begin{solution}{iq:s1:volatility-targeting-and-risk-managed-portfolios:1}
Si les rendements espérés ne croissent pas avec la volatilité, une période agitée comporte davantage de risque par unité de rendement espéré~; réduire l'exposition retire alors du risque
plus vite que le rendement, et le ratio s'améliore dans le temps.
\end{solution}

\begin{solution}{iq:s1:volatility-targeting-and-risk-managed-portfolios:2}
L'alpha provient d'une combinaison des portefeuilles pilotés et non pilotés, dont les pondérations sont ajustées sur l'ensemble de l'échantillon et ne sont pas disponibles en temps réel~;
hors échantillon, les pondérations de la combinaison sont instables et l'avantage disparaît en général.
\end{solution}

\begin{solution}{iq:s1:volatility-targeting-and-risk-managed-portfolios:3}
Fonds de \omterm{def:s1:trend-following:voltarget}{ciblage de volatilité}, parité de risque, ETF à effet de levier remis à zéro à la clôture, et dealers d'options en position courte de gamma~; il faut estimer
les actifs de chaque groupe, sa règle et sa fenêtre de prévision, et appliquer le mouvement du jour pour obtenir la transaction de demain.
\end{solution}

\begin{solution}{iq:s1:volatility-targeting-and-risk-managed-portfolios:4}
Elle réduit l'exposition aux périodes de forte volatilité et la gravité des pertes extrêmes~; elle ajoute de la rotation, vend après les baisses, accroît l'effet de levier dans les périodes calmes avant
les chocs, et lie le portefeuille à des flux très fréquentés.
\end{solution}

\begin{solution}{iq:s1:volatility-targeting-and-risk-managed-portfolios:5}
Collecter les cours de clôture, mettre à jour la prévision, calculer l'exposition cible avec les caps et les bandes, contrôler les limites et la liquidité, envoyer les ordres près de la clôture
avec une solution de repli si les marchés sont perturbés, et consigner chaque étape pour le rapprochement.
\end{solution}

\begin{solution}{iq:s1:volatility-targeting-and-risk-managed-portfolios:6}
Avec $w_t = c/\sigma_t$, le rendement vaut $c\mu/\sigma_t + c\varepsilon_t$, de moyenne $c\mu E[1/\sigma]$ et de variance $c^2(1 + \mu^2\operatorname{Var}(1/\sigma))$.
La position constante a pour ratio de Sharpe $\mu/\sqrt{E[\sigma^2]}$. Pour $\mu$ petit, le rapport des deux vaut $E[1/\sigma]\sqrt{E[\sigma^2]} \ge E[1/\sigma]E[\sigma]
\ge 1$ d'après l'inégalité de Jensen, avec égalité seulement lorsque $\sigma_t$ est constant.
\end{solution}
